\subsection{Concepto de administración de operaciones}
La administración de operaciones es el conjunto de actividades mediante las cuales una organización convierte insumos, como materiales, mano de obra, tecnología e información, en bienes o servicios que el cliente necesita. Esta función existe en toda empresa, sin importar si produce zapatos, ofrece servicios financieros o distribuye alimentos. Siempre hay un proceso que transforma algo en un resultado útil, y la tarea del administrador de operaciones es asegurarse de que ese proceso funcione bien \parencite{chase2014}.

Desde una perspectiva gerencial, la administración de operaciones no se reduce a supervisar una línea de producción. Implica tomar decisiones sobre cómo diseñar los procesos, qué tecnología utilizar, cómo organizar el espacio físico, cómo controlar la calidad y cómo coordinar a las personas. Estas decisiones tienen impacto directo en los costos, la calidad y el tiempo de entrega, tres factores que determinan en buena medida la competitividad de la empresa.

\subsection{Importancia para la empresa}
Una empresa que gestiona bien sus operaciones puede ofrecer productos o servicios de mejor calidad, a menor costo y con mayor rapidez que sus competidores. Esto se traduce directamente en clientes más satisfechos, menores desperdicios y mejores márgenes de ganancia. Por el contrario, cuando las operaciones son desorganizadas, los problemas se acumulan: retrasos en entregas, fallas en la calidad, costos elevados y pérdida de clientes \parencite{heizer2014}.

Un ejemplo sencillo puede ilustrar esto. Una empresa de manufactura que coordina bien su inventario de materias primas evita paros en la producción por falta de insumos. Al mismo tiempo, si no acumula exceso de materiales, reduce costos de almacenamiento. Este equilibrio, aunque parece simple, requiere decisiones constantes y bien fundamentadas. De la misma manera, un hospital que programa correctamente sus turnos de enfermeras y el uso de quirófanos puede atender a más pacientes con los mismos recursos.

La administración de operaciones también influye en la capacidad de innovar. Una empresa que conoce bien sus procesos puede identificar dónde mejorar, qué eliminar y qué automatizar. Esto la hace más ágil y mejor preparada para adaptarse a los cambios del mercado o la tecnología.

\subsection{Principales áreas de decisión}
Los gerentes de operaciones enfrentan decisiones en varias áreas a la vez. Podemos identificar diez áreas de decisión que abarcan desde el diseño del producto hasta el mantenimiento de equipos \parencite{heizer2014}. A continuación se describen las más relevantes desde un enfoque empresarial:

\textbf{Diseño del producto o servicio.} Antes de producir, la empresa debe decidir qué ofrece y cómo. Este diseño define las características del bien o servicio y establece las bases para todo el proceso productivo. Un producto mal diseñado desde el inicio genera problemas en producción, en calidad y en la satisfacción del cliente.

\textbf{Gestión de la calidad.} La calidad no es solo el resultado final; es un proceso continuo. La empresa debe definir estándares, establecer controles y asegurarse de que cada etapa del proceso contribuya al resultado esperado. En muchas organizaciones, se prefiere prevenir los errores antes de que ocurran, en lugar de corregirlos al final, lo que resulta mucho más costoso.

\textbf{Diseño del proceso.} Esta área responde a la pregunta de cómo se hará el trabajo. Puede ser un proceso manual o automatizado, continuo o por lotes, estandarizado o flexible. La elección depende del volumen de producción, el tipo de producto y los recursos disponibles. Un proceso bien diseñado mejora la productividad y reduce el desperdicio.

\textbf{Gestión del inventario y la cadena de suministro.} Las empresas necesitan insumos para producir, y esos insumos provienen de proveedores. Administrar bien esta cadena significa asegurarse de que los materiales lleguen a tiempo, en las cantidades correctas y al mejor costo posible. Un fallo en la cadena de suministro puede detener toda la producción. Por eso, muchas empresas establecen relaciones de largo plazo con proveedores confiables \parencite{krajewski2013}.

\textbf{Programación de la producción.} Esta área define cuándo se produce, en qué orden y con qué recursos. Una buena programación reduce tiempos de espera, evita cuellos de botella y asegura que los pedidos se entreguen puntualmente. Es especialmente importante en empresas con alta demanda variable o múltiples líneas de producción.

\subsection{Operaciones en empresas de manufactura y de servicios}
Aunque históricamente la administración de operaciones se asoció con la industria manufacturera, hoy se aplica con igual relevancia en empresas de servicios. La diferencia principal está en que los servicios suelen ser intangibles y se producen y consumen al mismo tiempo, lo que dificulta almacenarlos o corregirlos después de entregados.

Un banco, por ejemplo, gestiona operaciones cuando procesa solicitudes de crédito, atiende clientes en ventanilla o ejecuta transferencias. Un hotel administra operaciones cuando organiza la limpieza de habitaciones, la asignación de personal y el servicio de alimentos. En todos estos casos, los principios de la administración de operaciones siguen siendo los mismos: usar bien los recursos, servir con calidad y hacerlo a tiempo \parencite{chase2014}.

La diferencia de aplicación entre manufactura y servicios exige que los gerentes adapten sus herramientas. En los servicios, la experiencia del cliente forma parte del proceso mismo, por lo que factores como la amabilidad, la rapidez y la consistencia son parte de la calidad operativa.

\subsection{Productividad y eficiencia}
Uno de los indicadores más usados en la administración de operaciones es la productividad, entendida como la relación entre lo que se produce y los recursos que se usaron para producirlo. Una empresa es más productiva cuando obtiene más resultados con los mismos o menos recursos. Mejorar la productividad no significa trabajar más, sino trabajar mejor: con mejores procesos, mejor distribución del trabajo, mejores herramientas o mejor uso del tiempo \parencite{heizer2014}.

La eficiencia, en cambio, compara lo que se logró con lo que era posible lograr en condiciones ideales. Una empresa puede ser eficiente sin ser productiva si sus estándares son bajos, o puede aspirar a ser más productiva sin descuidar la calidad. Por eso, los gerentes deben buscar un balance entre cantidad producida, calidad entregada y recursos consumidos.

En la práctica, mejorar la productividad puede venir de simplificar pasos innecesarios, eliminar esperas, automatizar tareas repetitivas o capacitar mejor al personal. Estas mejoras requieren observar el proceso con detalle y tomar decisiones fundamentadas en datos.
